\documentclass[10pt]{article}
\usepackage[T1]{fontenc}
\usepackage[utf8]{inputenc}
\usepackage[french]{babel}
\usepackage{amsmath,amssymb,mathrsfs,mathpazo}
\usepackage{geometry}
\geometry{a4paper,margin=2cm}
\usepackage{graphicx,array,tabularx,tabularray,colortbl,multirow,multicol,booktabs,enumitem,pifont,ulem,fancybox,diagbox,cancel}
\usepackage{pstricks,pstricks-add,pst-plot,pst-tree,pst-node,pst-text,pst-func,pst-math,pst-bspline,pst-3dplot,pst-eucl,pst-solides3d}
\let\pscalculate\undefined
\usepackage{pst-bezier}
\usepackage{tikz,pgfplots,tkz-tab}
\usetikzlibrary{arrows,arrows.meta,patterns,calc,positioning,trees,decorations.markings,shapes.geometric,shapes.misc,shapes.arrows,intersections,decorations.pathmorphing,decorations.pathreplacing,backgrounds,shadows,babel}
\pgfplotsset{compat=1.18}
\usepgfplotslibrary{fillbetween}
\usepackage[np]{numprint}
\usepackage{listings,makecell,pas-tableur,twoopt}
\let\barre\undefined
\usepackage[pstricks]{tablvar}
\definecolor{darkgreen}{RGB}{0,100,0}
\providecommand{\red}{\color{red}}
\providecommand{\blue}{\color{blue}}
\providecommand{\green}{\color{green}}
\providecommand{\black}{\color{black}}
\providecommand{\white}{\color{white}}
\providecommand{\gray}{\color{gray}}
\providecommand{\cyan}{\color{cyan}}
\providecommand{\magenta}{\color{magenta}}
\providecommand{\yellow}{\color{yellow}}
\providecommand{\violet}{\color{violet}}
\providecommand{\orange}{\color{orange}}
\providecommand{\euro}{\texteuro}
\providecommand{\EUR}{\texteuro}
\providecommand{\ieme}{\textsuperscript{e}}
\providecommand{\degres}{\ensuremath{^\circ}}
\providecommand{\degree}{\ensuremath{^\circ}}
\providecommand{\vv}[1]{\overrightarrow{#1}}
\providecommand{\Coord}[1]{#1}
\providecommand{\decosix}{}
\providecommand{\starredbullet}{\ensuremath{\bullet}}
\pagestyle{empty}
\setlength{\parindent}{0pt}
\usepackage{adjustbox}\DeclareUnicodeCharacter{2236}{:}\DeclareUnicodeCharacter{2190}{\ensuremath{\leftarrow}}
\providecommand{\R}{\mathbb{R}}
\providecommand{\N}{\mathbb{N}}
\providecommand{\Z}{\mathbb{Z}}
\providecommand{\Q}{\mathbb{Q}}
\providecommand{\D}{\mathbb{D}}
\providecommand{\C}{\mathbb{C}}
\providecommand{\e}{{\mathrm{e}}}
\providecommand{\vect}[1]{\overrightarrow{#1}}
\providecommand{\vectt}[1]{\overrightarrow{\mathrm{#1}}}
\providecommand{\vvt}[1]{\overrightarrow{\mathrm{#1}}}
\providecommand{\np}[2][]{{#2~#1}}
\providecommand{\vv}[1]{{\overrightarrow{#1}}}
\renewcommand{\d}{{\mathrm{d}}}
\providecommand{\degres}{{^\circ}}
\providecommand{\degre}{{^\circ}}
\providecommand{\blue}{\color{blue}}
\providecommand{\red}{\color{red}}
\providecommand{\enskip}{\quad}
\providecommand{\llbracket}{\lbrack\!\lbrack}
\providecommand{\rrbracket}{\rbrack\!\rbrack}
\providecommand{\euro}{{\text{€}}}
\providecommand{\no}{{\text{n°}}}
\providecommand{\interval}[2]{{[#1;#2]}}
\providecommand{\intervFO}[2]{{[#1;#2[}}
\providecommand{\dotfill}{\ldots}

\providecommand{\ds}{\displaystyle}
\providecommand{\textding}[1]{#1}
\providecommand{\Oij}{\(\mathrm{O};\vec{\imath},\vec{\jmath}\)}
\providecommand{\Oijk}{\(\mathrm{O};\vec{\imath},\vec{\jmath},\vec{k}\)}
\providecommand{\Ouv}{\(\mathrm{O};\vec{u},\vec{v}\)}
\providecommand{\intervFF}[2]{\ensuremath{[#1;#2]}}
\renewcommand{\euro}{\ifmmode\text{\texteuro}\else\texteuro\fi}
\DeclareUnicodeCharacter{00B0}{\ifmmode^{\circ}\else\textdegree\fi}
\DeclareUnicodeCharacter{20AC}{\ifmmode\text{\texteuro}\else\texteuro\fi}
% Énoncés : Éducation nationale. Transcription et corrigés : APMEP, Denis Vergès et contributeurs. Découpage et classement : CoopMaths, Sébastien Lozano, Éric Elter et contributeurs. Adaptation de la présentation : Maths Explorer.
% CC BY-SA : https://coopmaths.fr/www/about/
% https://coopmaths.fr/alea/static/e3c/2020/tex/e3c_2020_00_sujet20_3_cor.tex
\begin{document}
\noindent\textit{APMEP / CoopMaths. CC BY-SA. Présentation adaptée par Maths Explorer.}\par\medskip
\subsection*{1.}
	
	Diminuer de 2 \% revient à multiplier par \(1 - 0,02 = 0,98\).
	
	\[
	120 \times 0,98 = 117,6 \, \text{kg} \quad \text{et} \quad 117,6 \times 0,98 = 115,248 \, \text{kg}.
	\]
	
	\subsection*{2.}
	\subsubsection*{a.}
	Diminuer de 2 \% revient à multiplier par \(1 - 0,02 = 0,98\).
	
	On a donc pour tout \(n \in \mathbb{N}\), \(a_{n+1} = 0,98a_n\) : la suite \((a_n)\) est une suite géométrique de premier terme \(a_0 = 120\) et de raison 0,98.
	
	\subsubsection*{b.}
	
	On sait alors que quel que soit le naturel \(n\),
	\[
	a_n = a_0 \times 0,98^n = 120 \times 0,98^n.
	\]
	
	\subsubsection*{c.}
	
	On rappelle que :
	
	Soit \((a_n)_{n \in \mathbb{N}}\) une suite géométrique de raison \(q\), \(q \neq 1\). La somme \(S\) de termes consécutifs est égale à
	\[
	S = u_1 + u_2 + \ldots + u_n = u_1 \times \dfrac{1 - q^n}{1 - q}.
	\]
	
	En 2020, la famille A produira :
	\[
	S_{2020} = a_1 + a_2 + \ldots + a_{12} = 120 \times 0,98 + 120 \times 0,98^2 + \ldots + 120 \times 0,98^{12}.
	\]
	
	Or
	\[
	0,98S_{2020} = 120 \times 0,98^2 + \ldots + 120 \times 0,98^{12} + 120 \times 0,98^{13}.
	\]
	
	Par différence entre les deux lignes précédentes :
	\[
	0,02S_{2020} = 120 \times 0,98 - 120 \times 0,98^{13},
	\]
	
	ou encore :
	\[
	0,02S_{2020} = 120 \times 0,98(1 - 0,98^{12}),
	\]
	
	et
	\[
	S_{2020} = \dfrac{120 \times 0,98(1 - 0,98^{12})}{0,02} = 49 \times 120 (1 - 0,98^{12}) \approx 1265,8.
	\]
	
	Donc \(S_{2020} \approx 1266 \, \text{kg}\) de déchets en 2020.
	
	\subsubsection*{d.}
	

	
	S(6) donne la somme \(a_1 + a_2 + \ldots + a_6\), soit la somme des déchets de la famille du premier semestre 2020.

\end{document}
